现在的任务显而易见 (也显然并不轻松). 必须将微分方程 $L(\bm{x},D)u=f$ 变换到 $\bm{y}$ 坐标中. 为便于这一工作, 引入如下记号: 对任意函数\[v:\Omega^{\prime}\to\mathbb{R},\]定义
\begin{equation}\label{eq:boundary-coordinate-transformed-function-domain}
\widetilde v:\Omega^{\prime\prime}\to\mathbb{R}
\end{equation}
使得
\begin{equation}\label{eq:boundary-coordinate-transformed-function-definition}
\widetilde v(\bm{y}):=v(\overline\Phi(\bm{y})).
\end{equation}
注意, 对任意函数 $v\in L^2(\Omega)$, 存在常数 $c_1$ 和 $c_2$, 使得
\begin{equation}\label{eq:boundary-coordinate-l2-norm-comparison}
c_1\|v\|_{L^2(\Omega^{\prime})}\leq\|\widetilde v\|_{L^2(\Omega^{\prime\prime})}\leq c_2\|u\|_{L^2(\Omega^{\prime})}.
\end{equation}


变量代换对偏微分算子的作用由以下引理描述.

\begin{Lemma}
\label{lem:boundary-coordinate-variational-transformation}
设 $u\in H^1(\Omega^{\prime})$ 在 $\partial\Omega\cap\partial\Omega^{\prime}$ 上满足 $u\equiv0$ (在迹的意义下), 且 $u$ 对所有 $v\in H_0^1(\Omega^{\prime})$ 满足变分方程
\begin{equation}\label{eq:boundary-coordinate-original-variational-equation}
B[v,u]=(f,v),
\end{equation}
则 $\widetilde u\in H^1(\Omega^{\prime\prime})$ 在 $\partial\Omega^{\prime\prime}\cap\{\bm{y}\mid y_n=0\}$ 上满足 $\widetilde u\equiv0$, 并且 $\widetilde u$ 对每个 $\widetilde v\in H_0^1(\Omega^{\prime\prime})$ 满足变分方程
\begin{equation}\label{eq:boundary-coordinate-transformed-variational-equation}
\widetilde B[\widetilde v,\widetilde u]=(\widetilde f,\widetilde v),
\end{equation}
这里
\begin{equation}\label{eq:boundary-coordinate-transformed-bilinear-form}
\begin{aligned}
\widetilde B[\widetilde v,\widetilde w]&:=-\sum_{k,l=1}^n\int_{\Omega^{\prime\prime}}\overline a_{kl}(\bm{y})\widetilde w_{y_l}\widetilde v_{y_k}\,\mathrm{d}\bm{y}\\
&\quad+\sum_{k=1}^n\int_{\Omega^{\prime\prime}}\overline b_k(\bm{y})\widetilde w_{y_k}\widetilde v\,\mathrm{d}\bm{y}\\
\mathit{quad}&+\int_{\Omega^{\prime\prime}}\overline c(\bm{y})\widetilde v\widetilde w\,\mathrm{d}\bm{y},
\end{aligned}
\end{equation}
其中
\begin{equation}\label{eq:boundary-coordinate-transformed-principal-coefficients}
\overline a_{kl}:=\sum_{i,j=1}^n\widetilde a_{ij}(\bm{y})\frac{\partial\Psi_k}{\partial x_i}(\Phi(\bm{y}))\frac{\partial\Psi_l}{\partial x_j}(\Phi(\bm{y})), 
\end{equation}

\begin{equation}\label{eq:boundary-coordinate-transformed-first-order-coefficients}
\overline b_k:=\sum_{i=1}^n\widetilde b_i(\bm{y})\frac{\partial\Psi_k}{\partial x_i}(\Phi(\bm{y})), 
\end{equation}

\begin{equation}\label{eq:boundary-coordinate-transformed-zero-order-coefficient}
\overline c(\bm{y}):=\widetilde c(\bm{y}).
\end{equation}

\end{Lemma}

证明使用标准技巧, 留给读者完成.
